% ============================================================
% CAPITOLO 4
% Casi europei e istituzionali di IA on-premise nella PA
% ============================================================
\chapter{Casi europei e istituzionali di IA on-premise nella Pubblica Amministrazione}
\label{chap:casi-europei}

% --- Introduzione al capitolo ---
Questo capitolo presenta una rassegna di casi europei e istituzionali verificati in cui sistemi di Intelligenza Artificiale sono stati adottati in modalità on-premise (o con componenti on-premise significative) all'interno di organizzazioni pubbliche. L'obiettivo è identificare best practice, modelli replicabili, fattori di successo e lezioni apprese, con particolare attenzione al contesto italiano ed europeo.


% ============================================================
\section{Metodologia di selezione dei casi}
\label{sec:metodologia-casi}

% --- Da completare ---
% Contenuti suggeriti:
% - Criteri di selezione: verificabilità, rilevanza, replicabilità
% - Fonti utilizzate: report istituzionali, pubblicazioni accademiche,
%   documenti ufficiali UE, OECD, JRC
% - Classificazione dei casi: per paese, per ambito applicativo,
%   per livello di maturità


% ============================================================
\section{Casi a livello europeo}
\label{sec:casi-europei}

% --- Da completare ---
% Contenuti suggeriti per ciascun caso:
% - Paese, ente, anno di avvio
% - Obiettivo e ambito applicativo
% - Architettura tecnologica (on-premise / hybrid)
% - Risultati ottenuti
% - Fonti verificate

\subsection{[Caso 1 — Da identificare]}
\label{subsec:caso-eu-1}
% ...

\subsection{[Caso 2 — Da identificare]}
\label{subsec:caso-eu-2}
% ...

\subsection{[Caso 3 — Da identificare]}
\label{subsec:caso-eu-3}
% ...


% ============================================================
\section{Casi a livello italiano}
\label{sec:casi-italiani}

% --- Da completare ---
% Contenuti suggeriti:
% - Casi di PA italiane con deployment on-premise di IA
% - Ministeri, Regioni, Atenei, enti locali
% - Progetti PNRR con componente IA on-premise

\subsection{[Caso italiano 1 — Da identificare]}
\label{subsec:caso-it-1}
% ...

\subsection{[Caso italiano 2 — Da identificare]}
\label{subsec:caso-it-2}
% ...

\subsection{[Caso italiano 3 — Da identificare]}
\label{subsec:caso-it-3}
% ...


% ============================================================
\section{Casi nel settore accademico}
\label{sec:casi-accademici}

% --- Da completare ---
% Contenuti suggeriti:
% - Atenei europei e italiani con infrastruttura IA on-premise
% - HPC e GPU cluster per ricerca e servizi
% - Esperienze di data center accademici

\subsection{[Caso accademico 1 — Da identificare]}
\label{subsec:caso-acc-1}
% ...

\subsection{[Caso accademico 2 — Da identificare]}
\label{subsec:caso-acc-2}
% ...


% ============================================================
\section{Analisi comparativa e best practice}
\label{sec:analisi-comparativa}

% --- Da completare ---
% Contenuti suggeriti:
% - Tabella comparativa dei casi (paese, ente, tecnologia, risultati)
% - Fattori di successo comuni
% - Criticità ricorrenti
% - Pattern architetturali emergenti (edge, federated, hybrid)
% - Best practice per l'adozione on-premise nella PA

% Esempio di tabella comparativa:
% \begin{table}[H]
% \centering
% \caption{Sintesi comparativa dei casi di IA on-premise nella PA}
% \label{tab:confronto-casi}
% \begin{tabularx}{\textwidth}{lXXXX}
% \toprule
% \textbf{Caso} & \textbf{Paese} & \textbf{Ambito} & \textbf{Tecnologia} & \textbf{Risultati} \\
% \midrule
% ... & ... & ... & ... & ... \\
% \bottomrule
% \end{tabularx}
% \end{table}


% ============================================================
\section{Lezioni apprese e raccomandazioni}
\label{sec:lezioni-apprese}

% --- Da completare ---
% Contenuti suggeriti:
% - Lezioni trasversali dai casi analizzati
% - Raccomandazioni per le PA italiane
% - Fattori abilitanti: governance, competenze, infrastrutture, budget
% - Implicazioni per il caso Parthenope (collegamento al Capitolo 3)
% - Prospettive di evoluzione
